\section{Supplementary Results and Research Leads}
\label{app:supplementary-register-02}

This appendix keeps potentially relevant results from the research
archive attached to a mathematical home without enlarging the main
theorem spine.  This register is published separately from the six reader manuscripts.  Inclusion here is deliberately permissive.  Each entry
states its evidence boundary; entries labelled as open, contextual, or
evidence-only are not used as theorems elsewhere in the paper.  Earlier
formulations known to be false are replaced by their current correction.
The computational supplement and its \texttt{COMPUTATION.md} index give
the available program-level artifact map; an entry here does not turn a
shared-code check into an independent verification.
\begingroup\raggedright\emergencystretch=3em

\subsection{Additional results and constructions}

\phantomsection
\label{supp-note-02-001}
\paragraph{After collision normalization, fixed-degree Keller counterexample...}
\emph{Status.} Proof offered; independent review is not recorded.

After collision normalization, fixed-degree Keller counterexample searches
become finite coefficient schemes; the archived derivation gives concrete
variable counts for \(D=4,5,6\) and recommends exact elimination or
modular-to-rational computation.

\phantomsection
\label{supp-note-02-002}
\paragraph{K\_\{2, 2\} is, after linear elimination, a twisted-cubic affine cone...}
\emph{Status.} Proof offered; independent review is not recorded.

K\_\{2, 2\} is, after linear elimination, a twisted-cubic affine cone: codimension two with three minimal quadratic generators, so it is not a complete intersection at the vertex.

\phantomsection
\label{supp-note-02-003}
\paragraph{An ordered-collision incidence construction makes the locus of...}
\emph{Status.} Proof offered; independent review is not recorded.

An ordered-collision incidence construction makes the locus of generic degree at least d open on the generically finite part; exact generic degree is locally closed and generic degree is lower semicontinuous in this setup.

\phantomsection
\label{supp-note-02-004}
\paragraph{The degree-four composite-pencil analysis gives a census/reduction...}
\emph{Status.} Recorded result or structural lead; it is not promoted beyond the evidence described here and in the project archive.

The degree-four composite-pencil analysis gives a census/reduction and restrictions on primitive-line and rank-one top forms.

\phantomsection
\label{supp-note-02-005}
\paragraph{If the quartic top form is rank one of pure-fourth-power...}
\emph{Status.} Proof offered; independent review is not recorded.

If the quartic top form is rank one of pure-fourth-power type H4=v times l\textasciicircum{}4, the Keller map is an automorphism.

\phantomsection
\label{supp-note-02-006}
\paragraph{For the triple-linear top form h=x\textasciicircum{}3 y, the generic charts...}
\emph{Status.} Archived chart-level research lead; no self-contained proof or
definition of the two exceptional divisors is offered here.

For the triple-linear top form \(h=x^3y\), the archived chart analysis
reports two residual exceptional divisors denoted \(\Delta_{23}\) and
\(\Psi\).  Those symbols are local to that archived elimination and are
not used by any theorem in this manuscript.

\phantomsection
\label{supp-note-02-007}
\paragraph{Every rank-one quartic map with top form v times l\textasciicircum{}3...}
\emph{Status.} Proof offered; independent review is not recorded.

Every rank-one quartic map in the treated orbit---namely one whose leading
map is \(H_4=v\ell^3m\), with \(v\ne0\) a fixed target vector and
\(\ell,m\) independent source linear forms---is an automorphism.

\phantomsection
\label{supp-note-02-008}
\paragraph{The corrected cubic-pencil reduction says that dP wedge dQ wedge...}
\emph{Status.} Proof offered; independent review is not recorded.

The corrected cubic-pencil reduction says that dP wedge dQ wedge dh=0 forces either a nontrivial fixed component/gcd, a pencil-cube case, or dependence on two linear forms.

\phantomsection
\label{supp-note-02-009}
\paragraph{The fixed-component Pl\"ucker boundary for
\(h=x^2y^2\) is eliminated.}
\emph{Status.} Exact computer-assisted proof offered; independent review is
not recorded.

For \(H_4=(0,0,x^2y^2)\) with cubic pencil generated by \(x^2y\) and
\(xy^2\), the complete high-syzygy solution and lower Keller equations force
the third column of the linear part to vanish.  This closes the earlier
unresolved fixed-component chart.  It is a special rank-one calculation,
not a target-span-two theorem.

\phantomsection
\label{supp-note-02-010}
\paragraph{If the projective image of the quartic leading form H4...}
\emph{Status.} Proof offered; independent review is not recorded.

If the projective image of the quartic leading form H4 is a curve of degree c, then H4 factors as G times a degree-c binary parametrization evaluated on a coprime source pencil (A, B), with the exact degree relation deg(G)+c deg(A)=4.

\emph{Credit.} Credit recorded to ChatGPT (AI-assisted derivation and drafting).

\noindent
For notes 011--017, take a target normalization
\(H_4=(P,Q,0)\).  The \emph{characteristic Jacobian derivation} is
\(\delta=(\nabla P\times\nabla Q)\cdot\nabla\), the \emph{leading field} is
\(\C(P,Q)\), and its \emph{homogeneous period} is the positive generator
of the degrees of nonzero homogeneous elements in the relative algebraic
closure of \(\C(P,Q)\) in \(\C(x,y,z)\).  If
\(F=H_4+H_3+H_2+H_1\), a \emph{normal layer} is the image of a lower
homogeneous term \(H_j\) in the one-dimensional target quotient normal to
the leading target span; \(N_m\) denotes the first nonzero such layer, of
degree \(m\).  These definitions make the archived statements precise,
but the register does not supply the missing global reductions.

\phantomsection
\label{supp-note-02-011}
\paragraph{For a quartic leading map of generic Jacobian rank two...}
\emph{Status.} Recorded structural lead under the definitions above; not used
as a theorem without the leading-field reduction in the project archive.

For a quartic leading map of generic Jacobian rank two, the kernel of its characteristic Jacobian derivation is the relative algebraic closure of the leading field; the scalar degrees in this field form qZ for a homogeneous period q dividing 4.

\emph{Credit.} Credit recorded to ChatGPT (AI-assisted derivation and drafting).

\phantomsection
\label{supp-note-02-012}
\paragraph{For a quartic Keller deformation of a leading curve, the...}
\emph{Status.} Recorded structural lead; the determinant-order and
period-resonance argument is not reproduced here.

For a quartic Keller deformation of a leading curve, the first nonzero normal coefficient below determinant order nine can occur only at an order divisible by the homogeneous period; in the primitive period-four case the first three normal coefficients vanish.

\emph{Credit.} Credit recorded to ChatGPT (AI-assisted derivation and drafting).

\phantomsection
\label{supp-note-02-013}
\paragraph{The normal-resonance theorem extends to every degree d in dimension...}
\emph{Status.} Research lead extrapolated from the quartic argument; no
self-contained all-degree proof is offered here.

The normal-resonance theorem extends to every degree d in dimension three, with determinant threshold 3(d-1); consequently no three-variable Keller map in any degree can have a period-primitive leading line.

\emph{Credit.} Credit recorded to ChatGPT (AI-assisted derivation and drafting).

\phantomsection
\label{supp-note-02-014}
\paragraph{For a quartic Keller map with leading line H4=(P, Q, 0), the...}
\emph{Status.} Recorded structural lead; the characteristic-field extension
argument is not reproduced here.

For a quartic Keller map with leading line H4=(P, Q, 0), the first nonzero normal layer Nm is algebraic but not rational over C(P, Q), and the characteristic-field extension has degree at least 4/gcd(4, m).

\emph{Credit.} Credit recorded to ChatGPT (AI-assisted derivation and drafting).

\phantomsection
\label{supp-note-02-015}
\paragraph{The first-normal analysis reduces quartic leading-line maps to...}
\emph{Status.} Recorded routing lead; compare the proved four-loci reduction
in the main text, but do not treat this broader sentence as a theorem.

The first-normal analysis reduces quartic leading-line maps to binary pencils, one quadratic-source composite type, fourth-power or quadratic-square special fibers according to the normal degree, and fixed components subject to explicit valuation divisibility.

\emph{Credit.} Credit recorded to ChatGPT (AI-assisted derivation and drafting).

\phantomsection
\label{supp-note-02-016}
\paragraph{In a quartic leading-line counterexample the first nonzero normal...}
\emph{Status.} Proof-offered research lead; the reader manuscript uses it
only as an explicitly qualified routing input in the remark entitled
``The vanishing-third-component leaf.''

In a quartic leading-line counterexample the first nonzero normal layer must have degree three: degree one makes the normal coordinate linear, while degree two makes it a quadratic coordinate and reduces the remaining map to an excluded low-degree plane Keller map.

\emph{Credit.} Credit recorded to ChatGPT (AI-assisted derivation and drafting).

\phantomsection
\label{supp-note-02-017}
\paragraph{In the primitive coprime quartic leading-line case, the surviving...}
\emph{Status.} Recorded structural lead; the claimed uniqueness of the
valuation partition is not proved in this register.

In the primitive coprime quartic leading-line case, the surviving normal-degree-three valuation pattern is uniquely the partition 2+1, reducing the leading data to H4=(ell\textasciicircum{}4, q2\textasciicircum{}2, 0) and (H3)3=ell q2 with ell not dividing q2.

\emph{Credit.} Credit recorded to ChatGPT (AI-assisted derivation and drafting).

\phantomsection
\label{supp-note-02-018}
\paragraph{No quartic Keller map has leading form z\textasciicircum{}2(x\textasciicircum{}2, xy, y\textasciicircum{}2), despite the...}
\emph{Status.} Exact or computer-assisted certificate offered; see the computation index for its scope and independence boundary.

No quartic Keller map has leading form z\textasciicircum{}2(x\textasciicircum{}2, xy, y\textasciicircum{}2), despite the period-two binary-cubic resonance allowed in its second normal layer.

\emph{Credit.} Credit recorded to ChatGPT (AI-assisted derivation and drafting).

\phantomsection
\label{supp-note-02-019}
\paragraph{No quartic Keller map has leading form (A\textasciicircum{}2, AB, B\textasciicircum{}2) for coprime...}
\emph{Status.} Proof offered; independent review is not recorded.

No quartic Keller map has leading form (A\textasciicircum{}2, AB, B\textasciicircum{}2) for coprime independent ternary quadratic forms A and B.

\emph{Credit.} Credit recorded to ChatGPT (AI-assisted derivation and drafting).

\phantomsection
\label{supp-note-02-020}
\paragraph{No quartic Keller map has either remaining period-one scalar-conic...}
\emph{Status.} Exact or computer-assisted certificate offered; see the computation index for its scope and independence boundary.

No quartic Keller map has either remaining period-one scalar-conic leading form xy(x\textasciicircum{}2, xy, y\textasciicircum{}2) or x\textasciicircum{}2(x\textasciicircum{}2, xy, y\textasciicircum{}2).

\emph{Credit.} Credit recorded to ChatGPT (AI-assisted derivation and drafting).

\phantomsection
\label{supp-note-02-021}
\paragraph{No quartic Keller map in three variables has a nondegenerate...}
\emph{Status.} Census consequence of the reader's reduced-conic theorem and
\hyperref[supp-note-02-018]{notes 018}--
\hyperref[supp-note-02-020]{020}; the 018 and 020 legs are
certificate-dependent, and independent review is not recorded.

No quartic Keller map in three variables has a nondegenerate conic as the projective image of its highest homogeneous part.

\emph{Credit.} Credit recorded to ChatGPT (AI-assisted derivation and drafting).

\phantomsection
\label{supp-note-02-022}
\paragraph{No quartic Keller map in three variables has a nondegenerate...}
\emph{Status.} Exact or computer-assisted certificate offered; see the computation index for its scope and independence boundary.

No quartic Keller map in three variables has a nondegenerate rational cubic as the projective image of its highest homogeneous part.

\emph{Credit.} Credit recorded to ChatGPT (AI-assisted derivation and drafting).

\noindent
For notes 023--026, let
\(h=(h_0,h_1,h_2)\in\C[x,y]_4^3\) be a basepoint-free proper plane-quartic
parametrization and factor
\[
h_x\times h_y=\gamma n
\]
with \(n\) primitive.  The polynomial \(\gamma\) is the
\emph{ramification factor}.  The \emph{tangent-syzygy type} is the pair of
degrees of a minimal homogeneous Hilbert--Burch basis for the syzygy module
of \(n\).  This is the notation used in the Program~2 reader appendix
``The quartic leading-curve frontier and ramification filtration.''

\phantomsection
\label{supp-note-02-023}
\paragraph{A proper basepoint-free rational-quartic leading image of a quartic...}
\emph{Status.} Proof offered; independent review is not recorded.

A proper basepoint-free rational-quartic leading image of a quartic Keller
map cannot be immersed or simply ramified; equivalently, its ramification
factor \(\gamma\) has degree at least two.

\emph{Credit.} Credit recorded to ChatGPT (AI-assisted derivation and drafting).

\phantomsection
\label{supp-note-02-024}
\paragraph{If a proper rational-quartic leading image has ramification degree...}
\emph{Status.} Proof offered; independent review is not recorded.

If a proper rational-quartic leading image has ramification degree two, then its tangent-syzygy module cannot have unbalanced Hilbert-Burch type (1, 3); only type (2, 2) can remain at ramification degree two.

\emph{Credit.} Credit recorded to ChatGPT (AI-assisted derivation and drafting).

\phantomsection
\label{supp-note-02-025}
\paragraph{Projective duality bounds the ramification degree of a proper...}
\emph{Status.} Proof offered; independent review is not recorded.

Projective duality bounds the ramification degree of a proper rational plane
quartic by three.  Before the continuation calculation, the earlier Keller
exclusions reduced the rational-quartic frontier to
\((\deg\gamma,\text{ tangent-syzygy type})=(2,(2,2))\) and
\((3,(1,2))\); both are now excluded by the rational-quartic frontier
theorem in the reader manuscript.

\emph{Credit.} Credit recorded to ChatGPT (AI-assisted derivation and drafting).

\phantomsection
\label{supp-note-02-026}
\paragraph{A proper rational quartic with ramification degree three is...}
\emph{Status.} Proof offered; independent review is not recorded.

A proper rational quartic with ramification degree three is projectively the tricuspidal quartic, because its dual is a rational nodal cubic; the displayed parametrization has ramification factor xy(x-y) and tangent-syzygy type (1, 2).

\emph{Credit.} Credit recorded to ChatGPT (AI-assisted derivation and drafting).

\subsection{Open problems and resolved research prompts}

\phantomsection
\label{supp-note-02-027}
\paragraph{Ordinary total degree \(7\) is not known to be minimal...}
\emph{Status.} Open problem or unresolved assertion.

Ordinary total degree \(7\) is not known to be minimal in dimension three.

\phantomsection
\label{supp-note-02-028}
\paragraph{The tricuspidal quartic leading-image prompt is resolved.}
\emph{Status.} Resolved by an exact computer-assisted proof; independent
review is not recorded.

The type-\((1,2)\) tangent-syzygy calculation excludes the tricuspidal
quartic as the projective leading image of a quartic Keller map, as recorded
in the rational-quartic frontier theorem in the reader manuscript.

\emph{Credit.} Credit recorded to ChatGPT (AI-assisted research proposal).

\phantomsection
\label{supp-note-02-029}
\paragraph{The balanced ramification-degree-two prompt is resolved.}
\emph{Status.} Resolved by an exact computer-assisted proof; independent
review is not recorded.

The two-root and double-root normal forms exclude the balanced
\((2,(2,2))\) tangent-syzygy stratum as a rational-quartic leading image,
as recorded in the rational-quartic frontier theorem in the reader
manuscript.

\emph{Credit.} Credit recorded to ChatGPT (AI-assisted research proposal).

\subsection{Context, corrections, and evidence boundaries}

\phantomsection
\label{supp-note-02-030}
\paragraph{For fixed n, D, normalized determinant-one polynomial maps form a...}
\emph{Status.} Reported context; not independently established here.

For fixed n, D, normalized determinant-one polynomial maps form a finite-type coefficient scheme K\_\{n, D\}; its tangent equation at F is the coefficientwise vanishing of tr((JF)\textasciicircum{}\{-1\} JH), equivalently the differential of det J.

\phantomsection
\label{supp-note-02-031}
\paragraph{Scaling contracts every bounded-degree normalized Keller component...}
\emph{Status.} Reported context; not independently established here.

Scaling contracts every bounded-degree normalized Keller component toward the identity; a component meeting the counterexample locus is generically counterexample, and invariant regular functions are constant along the contraction.

\phantomsection
\label{supp-note-02-032}
\paragraph{If D\_min is the minimum possible maximum ordinary total degree...}
\emph{Status.} Reported context; not independently established here.

If D\_min is the minimum possible maximum ordinary total degree of a coordinate of a complex polynomial Keller counterexample in dimension three, then 4 <= D\_min <= 7.

\emph{Credit.} Credit recorded to Angelo Vistoli (attributed contributor); Levent Alpöge (attributed contributor).

\endgroup
